\documentclass{article}
\usepackage{amsmath,amssymb,hyperref}
\title{A pending destination can become a new rolling-budget dependency}
\date{26 November 2026}
\begin{document}
\maketitle
\paragraph{Outcome and contract slice.}
At a headless civil-inspection simulation cut the whole queue can wait
on one pose, while two other jobs have the data needed to run.
The earlier report then experiences an actual destination rebinding.
A projected dispatcher using only present endpoint scope can lose the
earlier report even if it shares the \emph{real} ledger. A safe
conservative partition holds another report if its separate alias
and time summaries could overlap. Preserving their joint support
admits some of that work immediately, without losing either report.
An equally informed direct schedule search matches. This is a
source-replay discriminator, not independent I1 or I2 acceptance.

\paragraph{Pins and outputs.}
Use selected native broker dispatch, receive calls, and radio
accounting from osrf/subt revision
\texttt{083c6a8acb0c25aa1a35e34517e628cd3b52ae77}, with the pinned
broker harness and source blob manifest. The headless fixture has
valid endpoints, binary-exact simulation cuts, deterministic service
for the modeled path and a common byte cap. A real
delivery for the reported result is the selected destination
application's native callback. No actual RF propagation or physical
vehicle arrival is claimed. The same \texttt{TeamNode} radio states,
rolling debit queues and application objects are shared by all
partitions within an arm. No placeholder pose and no synthetic reset
or copied capacity are introduced.

The first pending report \(A\) currently names a recipient with a
missing pose; the ready report \(B\) currently uses a different
recipient. Report \(C\) is actually disjoint from all admitted
continuations. A separate scoped broker call can serve \(C\)
at time zero and \(B\) there if admitted. Stock first asks for
all team poses and returns with all three pending. The old duty
requires an actual callback at its authorized later dispatch cut
\(t_A\), in \(0.75,1,1.25\) seconds. At that cut the fixture makes
the missing pose available, and either leaves \(A\)'s topic bound
to its earlier physical recipient or installs the real replacement
endpoint binding to \(B\)'s recipient. The fixture does not infer the
replacement permission, binding, or timing from the missing pose.
The reported scopes are observed registrations at the respective
native calls. Endpoint authority, complete allowed future history
set, and scheduling opportunities are authored assumptions.
B's remaining opportunity, if sent late, is \(t_A+1.25\), at most
2.5 seconds, within its declared deadline. C's early service mark
is 0.25 seconds.

These mission requirements are deadline and budget constraints,
not a rule that the old logical queue order holds forever after
rebinding. If a stakeholder requires strict physical order
even after expiry of all budget interactions, that stricter law
can forbid a move that this experiment accepts. This matters
when comparing with the previous
\href{https://github.com/eternal-november/eos-research-programmes/pull/567}{projected
dependency queue}, which explicitly fixed precedence. Here the
reports have different application effects and compete only
through the documented native debit keys. There is one actual
attempt per report and all payloads individually fit the cap.
No lost report is silently retried or counted as done.

\paragraph{Joint replay law.}
Native rolling accounting retains the boundary at exactly one
second: it removes strictly older debit entries. Even if current
logical names are different, a possible source tuple \(w\) can
share the true receiver \(\rho(w)\), true sender \(\sigma\), or both.
For these two reports of sizes \(a,b\) and capacity \(L\), sending
\(B\) at time zero preserves the old duty in all admitted histories
exactly when
\[
 \forall w\in S,\quad
 \bigl((\rho(w)\lor\sigma)\land t_A(w)\le 1\bigr)
 \ \Longrightarrow\ a+b\le L.
\]
Proof: if one row is shared and still active, the prior callback
has consumed exactly \(b\) bytes on that row. Otherwise it adds
nothing to the row at the old cut. Both sender and actual receiver
must pass, so either one can be sufficient to block delivery;
neither may be reset by splitting the queue. Here path acceptance
is fixed. In another envelope a failed send may still charge TX
but not RX; a future-scope claim must carry those distinct effects,
as in the prior
\href{https://github.com/eternal-november/eos-research-programmes/pull/618}{eligibility
replay}. If capacity depends on the eventual actual transmitter,
the tested homogeneous \(L\) cannot be used without adaptation.

The ordinary window calculation extends to a nonnegative replay
ledger, if the remaining reliable deliveries, actual receivers
and their timestamp and same-cut order are fixed per whole world.
For each due operation count its own charge, all preceding due
charges still in that ledger window and all earlier admitted
prefix charges still in the window. Check against the consuming
row's actual cap. If all rows pass, induction yields every due
delivery; at a first failing operation the promised sequence
cannot all deliver. This does not handle a planner that can
reroute or cancel failures after the cut without adding those
alternatives to the source law. Nor does it account for a hidden
effect before the native callback.

\paragraph{Four compared decisions and sharp witness.}
All arms have the same common physical budget and manifest.
After the pose clears even the stock arm is allowed to replan
the held queue with the due old job first and B's valid late
slot. The four decisions are:
\begin{enumerate}
\item Stock whole-pose wait. It preserves eventual outputs in
the positive model, but cannot deliver the disjoint report by its
early mark.
\item Shared-ledger projection using current endpoint scope and
allowed times. It does not clone resources; it can still make
an unsafe early decision when the future alias appears.
\item A robust partition using the product of allowed alias and
allowed time lists. It retains the actual common ledger and
later allowances. This is safe but may insert histories not in
the declared joint source.
\item Joint replay with exact supported tuples. It is also
compared against direct enumeration of both send schedules on
the very same tuples, keeping the real debit ordering.
\end{enumerate}
At \(L=5,\ a=b=3\), with separate senders, a genuinely admitted
alias at \(0.75\) causes the snapshot arm's old report to fail.
An alias at exactly 1 second still does. For support consisting
of the nonalias outcome at \(0.75\) and the alias outcome
at \(1.25\), both real histories permit the early send.
The conservative product combines early time from the first
with alias from the second and must hold. On the selected alias
late execution, real native callbacks for all three arrive
without a split physical capacity. The correlated rule earns
earlier service only when the assumption excluding the early
alias has a valid source warrant before choosing \(B\). It does
not undo a real send if a newly allowed authority transition
invalidates that warrant. Such a source change needs prevention,
an alternative completion or explicit loss accounting.

\paragraph{Verification and limits of the comparison.}
Each of the four arms replayed 12,992 held-in cases. The snapshot
lost the old callback on 1,904; the robust factorized and joint
arms lost none. The joint arm admitted \(B\) at time zero in 238
more checked cases than the factorized safe arm; both served
\(C\) early in every case. The direct same-information schedule
check matched. Counts are a finite parameter sweep, not a
distribution or empirical uptime. The two-support comparison
is a scoped improvement against separate marginals, not a
claim that every factored scheduler would discard the pairing.
Selective architecture baselines such as
\href{https://arxiv.org/abs/2312.04190}{switchable action dependency graphs}
and \href{https://lillyclark.github.io/files/ACHORD_RA_L_IROS_.pdf}{ACHORD}
remain comparison targets; the specific source semantics and
opportunity timing used here are not end-to-end tests of either.
Lean checks the active row, marginal soundness and strict
witness. A real generator of transferable future-binding and
timing certificates, an independent selected composite workload,
whole-stack improvement and inventory dispositions remain open.
\end{document}
